\subsection{Projective Origin of Fluctuations}
  \label{subsec:projective-origin-of-fluctuations}

  Vacuum fluctuations arise from the non-injectivity of the projection~$\Pi$ in
  regimes where no stable localized excitations are present.
  In such regimes, a wide multiplicity of distinct relational configurations of~$\chi$
  project onto the same effective spacetime description.
  Since admissible observables are functions on $\mathcal{O} = \mathrm{Im}\,\Pi$,
  not on~$\Omega_\chi$~\cite{Beau2026m,Beau2026l}, intra-fibre differences are by
  definition invisible to observation: the residual structure within each
  fibre~$\Pi^{-1}(o)$ constitutes a degree of freedom that no admissible observable
  can resolve.
  Any two such co-projecting configurations $\omega_1, \omega_2 \in \Pi^{-1}(o)$
  share a common projected part---which yields the stable observable~$o$---and carry
  residual structural differences that are not erased by~$\Pi$.
  The projected manifestation of these residual differences constitutes the observable fluctuation.
  No intrinsic indeterminacy of the pre-geometric substrate is required: $\chi$ does not
  vary at the substrate level.
  Fluctuations are therefore an exclusive feature of the projected description, and their
  magnitude is controlled by the degree of non-injectivity of~$\Pi$, quantified by the
  effective projective entropy~$S_\Pi$.

  The spatiotemporal character of fluctuations---their variation in both space and
  time---follows directly from the same mechanism.
  Spatial variations arise from residual relational differences between co-projecting
  configurations that are localized in distinct regions of the effective description.
  Temporal variations arise because co-projecting configurations may differ in their
  position along the projective ordering parameter~$\lambda$: such ordering differences,
  once projected, are precisely what the effective description identifies as differences
  in operational time~$\tau$.
  The time-dependence of vacuum fluctuations is therefore not a signature of any
  infra-physical dynamics of~$\chi$, but the effective translation of projective ordering
  differences between co-projecting configurations.

  A natural question arises at this point: why do projection residuals \emph{appear random}
  rather than repeating identically under successive projections?
  If the substrate~$\chi$ were stationary and the projection~$\Pi$ were applied repeatedly
  to the same fibre configuration, the residuals would indeed repeat identically.
  The appearance of randomness is therefore not a property of~$\chi$ itself, but a
  consequence of the fact that the effective projection is not applied at a fixed relaxation
  stage.
  More precisely, the projection must be understood as implicitly $\lambda$-dependent:
  one effectively probes a family of maps~$\Pi_\lambda$.
  At each effective step~$U_n$, one probes not a fixed fibre $\Pi^{-1}(o)$, but a slice of
  the co-projecting set determined by the current value of the relaxation-ordering
  parameter~$\lambda$.
  Between two successive steps $U_n$ and $U_{n+1}$, $\lambda$ has advanced, and the set
  of substrate configurations that co-project onto the same observable~$o$ is therefore
  not identical.
  The advance of~$\lambda$ is monotonic but remains unobservable at the effective level,
  since $\lambda$-ordering is precisely what the effective description identifies as time.
  Each reprojection thus samples a \emph{different slice} of an evolving fibre, and the
  residual differences between these slices are registered as fluctuations in the effective
  description.
  In this sense, the apparent stochasticity of fluctuations is the observable signature of
  the irreversible advance of~$\lambda$, rather than an intrinsic indeterminacy of~$\chi$.
  Equivalently, if the projective ordering could be suspended while reprojections continued,
  the same residuals would repeat identically.
  Fluctuations therefore arise because the projection cannot be iterated at
  fixed~$\lambda$: time, understood as the projected image of the ordering induced by~$\Pi$,
  is the minimal and irreducible source of the observed noise.

  \begin{remark}[Fluctuations as $\lambda$-induced fibre drift]
    The mechanism identified above admits a natural geometric reading.
    The family $\{\Pi_\lambda\}_{\lambda \in \Lambda}$ defines a $\lambda$-parametrised bundle
    of co-projecting fibres over the observable space~$\mathcal{O}$.
    Fluctuations correspond to the transverse variation of successive fibre slices along the
    relaxation flow, i.e.\ to the drift of~$\Pi_\lambda^{-1}(o)$ as $\lambda$ advances.
    A rigorous formulation of this drift---as a connection on the fibre bundle, or as a
    Lie-derivative along the relaxation vector field---is left as an open structural direction.
  \end{remark}

  Two distinct mechanisms may contribute to the apparent randomisation of successive
  projection residuals.
  The first is \emph{fibre decorrelation}: even if the norm of the residual differences
  remains approximately constant, the relaxation dynamics~$F$ generically drives
  successive fibre slices into increasingly orthogonal directions within~$\Omega_\chi$,
  so that successive residuals become uncorrelated without necessarily growing in
  amplitude.
  This decorrelation is sufficient to account for the apparent randomness of fluctuations
  and does not require any energy growth at the substrate level.
  The second, potentially distinct mechanism is \emph{projective amplification}: if the
  restriction of~$F$ to the fibre admits unstable directions in~$\Omega_\chi$, residual
  norms may grow along the relaxation sequence until bounded by the Born--Infeld
  admissibility constraint~$A^{\mathrm{BI}}_{\max}$.
  Once this ceiling is reached, the projection enters a saturation regime and structural
  compensation mechanisms become operative; the full hierarchy of saturation responses
  is developed in
  Section~\ref{subsec:structural-interpretation-projective-thermodynamics}.
  Whether projective amplification is active alongside decorrelation, and under what
  conditions each mechanism dominates, remains an open structural question.

  The Born--Infeld bound plays a dual role in this picture.
  It caps the amplitude of projective amplification, providing a finite structural ceiling
  on fluctuation growth.
  Crucially, this ceiling is not imposed by hand but follows from the admissibility
  constraint on the projection~$\Pi$ itself.
  In standard quantum field theory, the absence of such a bound forces the introduction
  of ultraviolet regulators and renormalisation procedures to control divergent loop
  integrals.
  In the present framework, renormalisation is not required as an independent procedure:
  the Born--Infeld bound replaces it with a structurally motivated finite cutoff,
  grounded in the projective architecture of the effective description.
